\subsubsection*{1.2.5.}
$Y_{\bar s}$ 的每个几何点 $\bar x$ 定义 $Y\times_s\eta$ 的一个几何点 $(\bar x,\bar\eta)$, 其对应的纤维函子为
\begin{equation*}
F_\eta\longmapsto(F_{\bar\eta})_{\bar x}.
\end{equation*}
它还通过
\begin{equation*}
F_{(\bar x,\bar\eta)}=(F_{\bar\eta})_{\bar x}
\qquad\text{及}\qquad
F_{(\bar x,\bar s)}=(F_{\bar s})_{\bar x}.
\end{equation*}
定义 $Y\times_sS$ 的几何点 $(\bar x,\bar\eta)$ 和 $(\bar x,\bar s)$. 可以验证, 拓扑斯 $Y\times_s\eta$ 和 $Y\times_sS$ 的所有点都由这种方式得到. 我们不使用这一事实. 还可立即验证, 这些点集是保守的.

\subsubsection*{${}^{*}$注 1.2.6.}
1.2.5 的构造给出从 $Y\times_s\eta$（分别 $Y\times_sS$）的点所成范畴 $\operatorname{Point}(Y\times_s\eta)$（分别 $\operatorname{Point}(Y\times_sS)$）到 2-纤维积
\begin{equation*}
\operatorname{Point}(Y)\times_{\operatorname{Point}(s)}
\operatorname{Point}(\eta).
\end{equation*}
的等价. 这一等价也可由 $Y\times_s\eta$（分别 $Y\times_sS$）是拓扑斯的 2-纤维积这一事实推出.${}^{*}$

\subsubsection*{1.2.7.}
由 2-纤维积的函子性, 图 (1.2.4.1) 关于 $Y$ 和 $S$ 具有函子性. 特别地:
\begin{enumerate}[label=\alph*)]
\item $s$ 上的每个概形态射 $f:Y\longrightarrow Y'$ 都诱导从 $Y\times_sS$ 到 $Y'\times_sS$ 的拓扑斯态射. 有
\begin{equation}\tag{1.2.7.1}
f^*(F_{\bar s},F_{\bar\eta},\varphi)
=(f^*F_{\bar s},f^*F_{\bar\eta},f^*(\varphi)),
\end{equation}
而当 $f$ 拟紧时有
\begin{equation}\tag{1.2.7.2}
f_*(F_{\bar s},F_{\bar\eta},\varphi)
=(f_*F_{\bar s},f_*F_{\bar\eta},f_*(\varphi)).
\end{equation}
对 $f:Y\times_s\eta\longrightarrow Y'\times_s\eta$ 也有类似的公式.

\item Hensel 离散赋值环谱的每个满射态射 $f:S'\longrightarrow S$ 都诱导从 $Y_{s'}\times_{s'}S'$ 到 $Y\times_sS$ 的拓扑斯态射. 仍以 $f$ 表示自然态射 $Y_{s'}=Y\otimes_{k(s)}k(s')\longrightarrow Y$, 则有
\begin{equation}\tag{1.2.7.3}
f^*(F_{\bar s},F_{\bar\eta},\varphi)
=(f^*F_{\bar s},f^*F_{\bar\eta},f^*(\varphi)).
\end{equation}
当 $f$ 有限时, $f_*$ 可表示为诱导表示.
\end{enumerate}

\subsubsection*{1.2.8.}
设 $f:Y\longrightarrow Y'$ 为 $s$-概形间的局部有限型分离态射. 定义“具有固有支集的正像”函子
\begin{equation*}
f_!: (\text{带基点集合层, 底为 }Y\times_sS)
\longrightarrow
(\text{带基点集合层, 底为 }Y'\times_sS)
\end{equation*}
其公式为
\begin{equation}\tag{1.2.8.1}
f_!(F_{\bar s},F_{\bar\eta},\varphi)
=(f_!F_{\bar s},f_!F_{\bar\eta},f_!(\varphi)).
\end{equation}
当 $f$ 为局部闭浸入时, $f_!$ 称为补零延拓函子.

对 $Y\times_s\eta$ 可定义类似的函子.

\subsubsection*{1.2.9.}
设 $f:Y\hookrightarrow Y'$ 为 $s$-概形的局部闭浸入. 函子 $f_!$ 有右伴随 $f^!$, 即“支集在 $Y$ 中的截面”, 它满足
\begin{equation}\tag{1.2.9.1}
f^!(F_{\bar s},F_{\bar\eta},\varphi)
=(f^!F_{\bar s},f^!F_{\bar\eta},f^!(\varphi)).
\end{equation}
现在只假设 $f$ 拟有限. 从 $Y\times_sS$ 上阿贝尔层到 $Y'\times_sS$ 上阿贝尔层的函子 $f_!$ 仍有由 (1.2.9.1) 给出的右伴随 $f^!$（参见 SGA~4, XVIII, 3.1.8）. 若 $f$ 平展, 则 $f^!=f^*$.

对 $Y\times_s\eta$ 可定义类似的函子.
